\sectionkicker{Theme synthesis}
\subsection*{横断比較 — Prompt Injectionはどの境界を越えるのか}
\addcontentsline{toc}{subsection}{横断比較 — Prompt Injectionはどの境界を越えるのか}
Agentがweb、tool、fileへ接続されると、prompt injectionは単なる誤応答ではなく情報flowとexternal actionの問題になる。1月の資料は、具体的なURL exfiltration対策と広い攻撃taxonomyを異なるEvidence classとして示した。
\medskip
\begin{center}
\small
\renewcommand{\arraystretch}{1.16}
\begin{tabularx}{\linewidth}{@{}p{0.20\linewidth}X@{}}
\toprule
観察軸 & 1月の一次資料・論文から読めること \\
\midrule
脅威の具体形 & OpenAIの1月28日の技術解説は、prompt-injected contentがAgentにURL取得を促し、会話由来の秘密をURLへ載せて外部送出し得るURL-based data exfiltrationを扱う。 \cite{src-0d2abbb0fa62} \\
data flow & このthreat modelでは秘密情報そのものが通常のresponse本文ではなく、Agentが取得するURLの一部として外部へ運ばれる。tool actionを持つAgentではoutput filteringだけでは捉えにくい情報flowになる。 \cite{src-0d2abbb0fa62} \\
automatic retrieval境界 & 説明されたsafeguardは、独立にpublicと確認されたURLだけを自動取得の候補とし、それ以外はfetchを遮断するか明示的user actionへ戻す。autonomyとconfirmationの境界をURL provenanceで切る設計である。 \cite{src-0d2abbb0fa62} \\
defense-in-depth & Safe URLはURL-based exfiltrationへの一層であり、任意のweb contentが安全になったことやprompt injection一般の解決を意味しない。security propertyもここでは独立penetration testしていない。 \cite{src-0d2abbb0fa62} \\
attack ecosystem & Prompt Injection SoKはagentic coding assistantだけでなくskills、tools、protocol ecosystemにまたがる攻撃を整理する。入力prompt一箇所ではなく、Agentが接続する周辺surfaceまで脅威モデルを広げている。 \cite{src-e9010a7fe30b} \\
taxonomy & 同SoKは3次元のattack taxonomyを提案し、異なる攻撃条件を整理する。ここで重要なのは単一attack success rateではなく、どのsurface・条件・mechanismを比較しているかを分離することである。 \cite{src-e9010a7fe30b} \\
literature scale & abstractが報告する78研究、42攻撃手法、18防御機構は著者によるliterature synthesisの集計であり、新規controlled benchmarkの再現値ではない。異質な研究を一つの実験結果へ潰さない。 \cite{src-e9010a7fe30b} \\
Evidence boundary & OpenAI資料は特定実装のdefense-in-depth、SoKは文献全体のtaxonomyという異なるEvidenceである。前者の実装claimと後者のsurvey集計を同じ強さの検証結果として扱わず、SoKは1月時点のarXiv v1へ固定する。 \cite{src-0d2abbb0fa62,src-e9010a7fe30b} \\
\bottomrule
\end{tabularx}
\renewcommand{\arraystretch}{1.0}
\end{center}
\normalsize
\begin{claimboundary}[この比較の境界]
この表は本号で既に検証・参照している一次資料と論文を横断比較の形へ再配置したものである。vendor / project / author が報告した性能値や優位性は独立再現済みの一般則へ変換せず、本文とSource Notesの帰属境界をそのまま維持する。
\end{claimboundary}
